\section{Lean formalization and the ProveIt integration path}\label{sec:lean}

\subsection{Three layers which must remain distinct}

The accompanying \texttt{lean/Propositional.lean} defines a readable
inductive formula type with falsehood, natural-number atom names, conjunction,
disjunction and implication. It defines Boolean evaluation, proposition-valued
interpretation, substitution, size, atom occurrences, and a mode-indexed
natural-deduction type. Its candidate proofs establish substitution laws and
proof soundness, with a decision assumption required for classical axioms.
The file does not import a classical axiom in order to prove intuitionistic
soundness. It is a \emph{specification-level} development in Lean's metatheory.

The \emph{numerical layer} represents all the data by the codes fixed in
Section~\ref{sec:base}, with primitive recursive bit checkers. This layer
must connect the readable inductive specifications to finite arrays and
bounded recursive algorithms. A total Lean function by itself is not
automatically a primitive recursive function, so the construction from
\texttt{PR.zero}, \texttt{PR.succ}, \texttt{PR.proj}, \texttt{PR.comp}, and
\texttt{PR.prec} must be retained.

The \emph{object-derivation layer} proves the numerical specifications in
the embedded arithmetic. Its final target is
\[
 \texttt{Prf PRAAx [] phi},\qquad \texttt{phi : Fm PTm 0}.
\]
This is a derivation of a specified closed sentence of iPRA. It is stronger
evidence for the requested formalization than an ambient Lean theorem
about the standard natural numbers. The present package supplies the
human proofs of those object-theory results and a detailed implementation
route, but not their completed object-level Lean derivations.

\subsection{Existing interfaces actually inspected}

The inspected \texttt{VladimirReshetnikov/ProveIt} snapshot is:
\begin{center}\texttt{803d951d15834da5571e26d1f39988d921b29760}\end{center}
The root README distinguishes
kernel-checked libraries from delivered research packets; we preserve that
distinction. We inspected the actual source of the following interfaces.
The full paths, blob identifiers, and roles are recorded in the source
audit included with the archive.

\begin{center}\small
\begin{tabularx}{\textwidth}{@{}>{\raggedright\arraybackslash}p{.29\textwidth}>{\raggedright\arraybackslash}X@{}}
\toprule
Module (abbreviated) & Verified interface and intended use\\
\midrule
\texttt{ConciseIPRA.PR} & Arity-indexed primitive recursive function syntax,
terms \texttt{PTm}, and the exact axiom family \texttt{PRAAx}.\\
\texttt{ConciseIPRA.Deduction} & Natural-deduction constructors and
\texttt{Prf.weaken}, \texttt{Prf.subst}, \texttt{Prf.trans}.\\
\texttt{Core.Pairing} & \texttt{pairT}, projections with their inverse
equations, and simultaneous primitive recursion.\\
\texttt{Core.Search} & \texttt{search\_spec},
\texttt{search\_found}, and \texttt{search\_minimal} for explicitly
bounded least-zero search.\\
\texttt{BooleanBits} & Bit closure and interpretation of negation,
conjunction, disjunction and Boolean implication.\\
\texttt{BoundedReflection} & Bounded universal/existential aggregators,
witness extraction, and bounded choice.\\
\texttt{BoundedSyntaxInduction} & Induction on a fixed bounded formula
through its actual characteristic-bit compilation.\\
\texttt{BoundedFixpoint} & Ascending finite-closure termination after
a proved finite rank bound; its hypotheses remain obligations for our
specific operator.\\
\bottomrule
\end{tabularx}\end{center}

The four foundation entries have prefix
\texttt{IPRADerivations.Foundations.}; the two entries marked
\texttt{Core} have prefix \texttt{IPRADerivations.}.
We use the source declarations as integration points; we did not rebuild
the upstream repository in this session and do not assert a new independent
verification of all its transitive dependencies.

\subsection{Exact numerical endpoint sentences}

Let $\mathrm{Fm}(a)$ be the formula-validity bit,
$\mathrm{Taut}(a)$ the truth-table product, $\mathrm{Prf}_C(d,a)$ the
empty-context classical proof-checker bit, and $\mathrm{makeC}(a)$ the
proof constructor, returning a default on nontautologies. The classical
completeness endpoint is the closed sentence
\[
 \forall a\,[\mathrm{Fm}(a)=1\wedge\mathrm{Taut}(a)=1
                   \to\mathrm{Prf}_C(\mathrm{makeC}(a),a)=1].
\]
It has only a quantifier-free matrix. Boolean soundness is a separate
closed endpoint, universally in $d,a$, and all required validity parameters.
To cover finite contexts, code the pair $(\Gamma,A)$ as a single input and
use the same construction.

For intuitionistic completeness, define $\mathrm{Prf}_I(d,a)$ and
$\mathrm{CModel}(c,a)$. The latter checks a nonempty finite partial order,
upward-closed valuations on the input atoms, a designated world, and
failure of $a$ there. The two output functions are $\mathrm{makeI}(a)$
and $\mathrm{model}(a)$, with defaults on their unused branches. The target is
\[
 \forall a\,[\mathrm{Fm}(a)=1\to
   (\mathrm{Prf}_I(\mathrm{makeI}(a),a)=1
       \vee\mathrm{CModel}(\mathrm{model}(a),a)=1)].
\]
The two displayed sentences are represented literally as closed
\texttt{Fm PTm 0} definitions in \texttt{ObjectExamples.lean}, with the
required primitive recursive symbols as parameters. The definitions are
not asserted as axioms or advertised as finished theorems. To finish
them, instantiate those parameters by the algorithms in this paper and
prove the stated checker contracts in \texttt{Prf PRAAx}.

\subsection{Definitions and proof order for completion}

The full plan in \texttt{lean/FormalizationPlan.md} gives the following
concrete work order.
\begin{enumerate}
\item \textbf{Coding.} Implement pairing, partial inverses with validity
tests, finite sequences, base-digit arrays, formula tags and decoders,
postorder/subcode iteration, duplicate-free supports and subformula lists,
and the proof/tree/DAG codes. Prove constructor disjointness, smaller
child codes and encoding/decoding round trips. A universal evaluator for
arbitrary primitive recursive program codes is neither required nor
available as a primitive recursive shortcut.
\item \textbf{Evaluation.} Construct the Boolean evaluator and finite
valuation enumeration. Prove its bit property, support locality and
substitution by bounded induction on a subcode table. Define row products,
truth tables, SAT, TAUT, and finite consequence; use the existing bounded
search interface to extract the least failed or successful row.
\item \textbf{Proof checking.} Define each natural-deduction rule tag
and its exact bounded side tests. Reject forward references and incorrect
conclusions or contexts. Prove primitive recursive weakening, substitution
and proof cut. Derive Boolean soundness using the invariant that every
earlier proof node is valid on every valuation row of a covering support.
\item \textbf{Classical completeness.} Implement the signed-evaluation
constructor and verify its two output modes. Implement the full binary
literal-table merge. Its invariant is the conjunction of checker bits at
the currently constructed frontier. This gives the classical endpoint,
after which many equivalence proof generators can simply reuse it.
\item \textbf{Intuitionistic completeness.} Define the $m2^m$-entry
sequent table, twelve-rule update, first-appearance records, $N$-fold
iteration, greedy prime extension, canonical-world predicate, model
checker and forcing evaluator. Instantiate the bounded-fixed-point
theorem, prove all five greedy invariants, and then prove the finite
truth lemma. Combine them into the second endpoint. Keep the explicitly
derived cut rule distinct from an unproved assumption of cut elimination.
\item \textbf{Other logical metatheorems.} Implement Hilbert line
translation, tableau search and LK reconstruction, negative translation,
and the stable-case proof merger. For disjunction extraction implement
generated cones and two-model gluing and verify the Harrop induction.
\item \textbf{Boolean algorithms.} Implement ordered normal forms,
ANF coefficient masks, Tseitin gate constraints, resolution annotations,
assignment restriction, DPLL fuel, Horn saturation, finite reachability
and SCC ordering, forgetting and definition extraction. The invariants
are the local equivalences and bounded model correspondences proved in
the relevant sections. Preserve the distinction between equivalence in
the old variables and equisatisfiability with new gate variables.
\item \textbf{Complexity reductions.} Implement the finite local machine
transition table, one-hot cell indicators, circuit unrolling and circuit
Tseitin encoding. Prove the bounded time/cell invariant and explicit
polynomial string-output bounds for each fixed verifier presentation.
\end{enumerate}

\texttt{ObjectExamples.lean} additionally supplies eight short object-level
proof candidates: identity, conjunction and disjunction commutativity,
implication composition, double-negation introduction, double-negated
excluded middle, explosion, and a closed universally quantified identity
for a supplied unary primitive recursive predicate. Their constructor
expressions use the inspected repository interface. They illustrate the
required endpoint discipline without pretending to establish the large
algorithmic theorems already.

\subsection{Executable evidence and its exact limits}

The Python reference program uses immutable formula trees, finite sets
of assumptions, finite proof DAGs, and bit tables. These are readable
representations of the specified finite objects. The accompanying README
explains the correspondence with numerical coding and the use of in-place
saturation scans, whose soundness and least-fixed-point interpretation
agree with the staged mathematical construction.

The main APIs return complete witnesses:
\begin{itemize}
\item \texttt{truth\_table\_prove(context,goal)} returns a constructed
classical proof DAG or a Boolean countervaluation.
\item \texttt{decide\_ipc(context,goal)} returns a pruned intuitionistic
proof DAG or a finite rooted countermodel with designated world zero.
\item Result checkers verify the exact requested context and conclusion;
JSON conversion permits certificate export and replay independently of
the producing function.
\end{itemize}
The demonstration file includes intuitionistic identity, intuitionistic
failure of excluded middle, a classical excluded-middle proof, and the
source article's explicit countervaluation.

The executed validation covers all 516 formulas with at most five
primitive nodes over two atoms and falsehood; 1,032 canonical-normal-form
equivalences; 51,600 substitution instances; 9,258 Tseitin extension
assignments; all 512 normalized two-variable clause sets; and 830 valid
entailments with interpolants. For 551 formula/context sequents, the
exporters produced 142 intuitionistic proof DAGs, 409 rooted countermodels,
150 classical proof DAGs and 401 Boolean countervaluations. All 1,102
serialized result round trips were replayed. Negative proof-checker
regressions cover wrong contexts, malformed discharge, incorrect targets,
self/forward/negative references, and classical axioms in intuitionistic
mode. Empty Kripke frames are rejected.

These checks exercise real certificates and the finite canonical truth
lemma. They are \emph{not} formal Lean proofs, an exhaustive test of
unbounded input sizes, or machine verification of the report's arithmetic
formalization. The report is a developed non-formalized proof package;
the Lean sources and object-level numerical endpoints remain the clear
next formalization task.
